\documentclass{report}

\usepackage{amsmath,amssymb,amsthm}
\usepackage{hyperref}
\usepackage[capitalize]{cleveref}

% Theorem environments
\newtheorem{theorem}{Theorem}
\newtheorem{definition}{Definition}
\newtheorem{axiom}{Axiom}
\newtheorem{corollary}{Corollary}
\newtheorem{remark}{Remark}

% Custom macros for leanblueprint
\providecommand{\home}[1]{}
\providecommand{\github}[1]{}
\providecommand{\dochome}[1]{}
\providecommand{\lean}[1]{}
\providecommand{\leanok}{}
\providecommand{\uses}[1]{}

\title{Structural Arithmetic and Continuous Approximation over the ISAR Kernel}
\author{Fabian Schindler}
\date{\today}

\begin{document}
\maketitle

\begin{abstract}
We bridge the discrete combinatory logic of the ISAR substrate with continuous analysis and ML theory. We define a structural quantity model and matrix representation of carriers, showing that matrix models form a terminal matrix kernel. We construct a pseudo-metric on address spaces, proving that its metric completion has a universal approximation property for continuous functions. This establishes a continuous representation limit for physical systems embedded in the ISAR substrate.
\end{abstract}

\chapter{Introduction and Foundational Philosophy}

Modern software engineering, programming language design, and formal verification are plagued by a profound fragmentation. Disjoint computational models—such as the untyped lambda calculus, term rewriting systems (TRS) and e-graphs, compiler intermediate representations (IR), arithmetic and quantity calculi, and virtual machine bytecodes—are treated as separate worlds. They occupy distinct namespaces, employ incompatible metatheories, and rely on ad-hoc, pairwise compiler pipelines that leak implementation details and fail to preserve observational boundaries.

This thesis argues that this fragmentation is unnecessary. We show that a broad class of closed formal systems can be represented through a single, quotient-mediated semantic kernel without loss of operational behavior. This kernel is not a new ``universal language,'' but rather a view-independent substrate: different programming languages and formalisms become different \textit{decoders} over the same invariant structure.

\section{Primal Axioms and Ontological Closure}

The ISAR (Invariant Scalable Fractal Addressable Representations) framework starts from the minimal assumptions of existence and distinction, leading to a mathematically closed computational substrate.

\begin{axiom}[Existence]
\leanok
The universe of discourse is non-empty: $\Omega \neq \emptyset$.
\end{axiom}

\begin{axiom}[Identity]
\leanok
The initial state is a singular, indistinguishable point (the Void, denoted $\emptyset$). Distinction begins when this Void is partitioned.
\end{axiom}

\begin{axiom}[Self-Description]
\leanok
A system exists only if it can describe its own existence. This self-description requirement forces the transition from a singleton to a pair: $P = (\emptyset, \emptyset)$.
\end{axiom}

\begin{axiom}[Asymmetry]
\leanok
The act of description creates a fundamental distinction between the \textit{Observer} and the \textit{Observed}. In set-theoretic terms, this corresponds to the Kuratowski pairing:
$$(x, y) = \{\{x\}, \{x, y\}\}$$
This asymmetry, denoted $\Delta$, serves as the operational fuel driving directed rewriting steps and computations.
\end{axiom}

\section{The 4-Carrier Minimality Proof by Exclusion}

The computational substrate is spanned by four primitive carriers: Identity ($I$), Rewrite ($R$), Adjacency ($A$), and State ($S$). We prove that exactly four independent carriers are necessary and sufficient to form an admissible computational substrate:

\begin{enumerate}
    \item \textbf{Identity ($C_I$)}: Represents the invariant operational equivalence. Without $C_I$, the substrate lacks a notion of stable normal forms or equivalence-preserving transformations, causing all state transitions to dissolve into noise.
    \item \textbf{Rewrite ($C_R$)}: Represents the directed, non-invertible selection/choice. Without $C_R$, transitions are symmetric, leading to static, reversible systems incapable of representing one-way computational steps.
    \item \textbf{State ($C_S$)}: Represents repeatable observables and state-space closure. Under self-application, the system must remain closed. Without $C_S$ to copy and distribute terms over contexts, the system collapses into a strictly linear, non-duplicating fragment that is sub-universal (incapable of simulating arbitrary Turing-complete computation).
    \item \textbf{Adjacency ($C_A$)}: Represents causal connectivity and topological interaction. It maps the relations between parameters. Without $C_A$, the carriers $C_I, C_R, C_S$ cannot be composed into ordered causal chains; the system collapses to an unordered bag of static values.
\end{enumerate}

No carrier can serve dual roles without violating single-carrier localization contexts:
\begin{itemize}
    \item If $C_I = C_S$, there is no proper quotient structure (observational equivalence collapses to identity).
    \item If $C_R = C_A$, causality becomes mutable, leading to non-deterministic, time-dependent topological rewrites that violate confluence.
    \item If $C_I = C_R$, the invariant becomes strategy-dependent, causing evaluation order to determine the semantic outcome.
\end{itemize}
Thus, the minimal size of the carrier set $C$ is exactly 4.

\section{Sizing Hierarchies: The Genesis of Nibbles and Bytes}

The token hierarchy of digital computing (bits, nibbles, bytes) is typically justified by legacy hardware constraints (e.g., ASCII encoding, memory alignment on the IBM System/360 in 1964). In the ISAR framework, this hierarchy is derived constructively from first principles of directed relation mapping:

\begin{itemize}
    \item \textbf{1 Bit}: Distinguishes existence from absence (atom vs. void).
    \item \textbf{2 Bits}: Identifies one of the four active roles ($C_I, C_R, C_A, C_S$) in the minimal carrier set.
    \item \textbf{4 Bits (One Nibble)}: Represents a directed pair of names (an input pointer to a carrier role and an output pointer from a carrier role). This is the minimum structure needed to express \textit{adjacency} (a relation between two roles).
    \item \textbf{8 Bits (One Byte)}: Represents a pair of pairs—a source context and a target context—which is the minimal faithful representation of a \textit{directed rewrite step} (a morphism between two carrier contexts).
\end{itemize}
Thus, early hardware engineers arrived at the 8-bit byte because it is the minimal structural representation of an asymmetric rewrite step ($\Delta$) with a parity check.

\section{Occam and the Description-Length Functional}

For any substrate $M$, we define the description-length marginal likelihood:
$$ p(D \mid M) = \int p(D \mid \theta, M) p(\theta \mid M) d\theta $$
where $\theta$ represents the parameter space (matrices and rules) and $D$ represents a reference computation. Given the 4-carrier limit, any substrate with $n > 4$ carriers introduces redundant gauge flexibility that increases description length. The terminal ISAR kernel minimizes this description length, acting as the unique Minimum Description Length (MDL) substrate.
\chapter{Continuous Approximation and Representation Space}

In this chapter, we bridge the discrete combinatory logic of the ISAR substrate with continuous analysis and matrix spaces. We formalize dimensional physical quantities, matrix representations, expressive completeness, the tensor semantics, and the continuous universal approximation theorem showing that every continuous function can be uniquely represented as a limit address.

\section{Physical Quantities and Matrix Representations}

We define dimensional quantities and 4x4 matrix algebra.

\begin{definition}[Quantities and Uncertainty]
\lean{ISAR.Quantity}
A quantity represents a physical parameter with:
\begin{enumerate}
    \item A value (rational scalar).
    \item A physical dimension (exponent exponents).
    \item An epistemic uncertainty expression \texttt{Uncertainty} and \texttt{Correlation}.
\end{enumerate}
Covariance propagation propagates exact metric-epistemic covariance equations under addition and multiplication.
\end{definition}

\begin{definition}[Matrix4 Carriers]
\lean{ISAR.Matrix4}
The concrete integer 4x4 matrix representation used for basis combinator signatures. The primitive matrices are:
\begin{gather*}
I_1 = \begin{pmatrix} 1 & 0 & 0 & 0 \\ 0 & 0 & 0 & 0 \\ 0 & 0 & 1 & 0 \\ 0 & 0 & 0 & 0 \end{pmatrix}, \quad
R_1 = \begin{pmatrix} 1 & 0 & 0 & 0 \\ 0 & 0 & 0 & 0 \\ 0 & 1 & 0 & 0 \\ 0 & 0 & 1 & 0 \end{pmatrix}, \\
A_1 = \begin{pmatrix} 0 & 0 & 0 & 0 \\ 1 & 0 & 0 & 0 \\ 0 & 1 & 0 & 0 \\ 0 & 0 & 0 & 0 \end{pmatrix}, \quad
S_1 = \begin{pmatrix} 1 & 1 & 0 & 0 \\ 0 & 1 & 0 & 0 \\ 0 & 0 & 1 & 0 \\ 0 & 0 & 0 & 1 \end{pmatrix}
\end{gather*}
\end{definition}

\begin{theorem}[Idempotency and Nilpotency]
\lean{ISAR.I1_idempotent, ISAR.K1_nilpotent}
\uses{ISAR.Matrix4}
The identity matrix $I_1$ is idempotent:
$$I_1^2 = I_1$$
The core rewrite operator $K_1 = I_1 \cdot R_1 \cdot A_1 \cdot S_1$ is nilpotent:
$$K_1^2 = \mathbf{0}$$
\end{theorem}

\begin{theorem}[Gauge Similarity]
\lean{ISAR.K1_K2_gauge_equiv}
\uses{ISAR.Matrix4}
The two alternative matrix representations $K_1$ and $K_2$ are conjugate (similar) via the lower-triangular gauge transformation matrix $P$:
$$P \cdot K_1 \cdot P^{-1} = K_2$$
proving that representation details are gauge shadows, while the terminal quotient state is invariant.
\end{theorem}

\begin{definition}[Matrix View Map]
\lean{ISAR.term_signature_val}
\uses{ISAR.Matrix4, ISAR.ITerm}
The structural homomorphism mapping each ISAR term to its concrete 4x4 integer matrix signature:
$$\text{term\_signature\_val} : \text{ITerm} \to \text{Matrix4}$$
\end{definition}

\begin{theorem}[Expressive Completeness]
\lean{ISAR.isk_expressive_completeness, ISAR.basis_expressive_completeness}
\uses{ISAR.term_signature_val}
Every matrix generated by the basis combinator algebra is the image of some term:
$$\forall M \in \text{ISKAlgebra}, \quad \exists t \in \text{ISKSubtype}, \quad \text{term\_signature\_val } t = M$$
\end{theorem}

\section{The Categorical Matrix Bridge}

We package the matrix representations as an admissible semantic kernel.

\begin{definition}[Matrix Kernel]
\lean{ISAR.MatrixKernel}
\uses{ISAR.Kernel, ISAR.term_signature_val}
The semantic kernel MatrixKernel whose carrier is \texttt{Matrix4} and whose view map is \texttt{term\_signature\_val}.
\end{definition}

\begin{theorem}[Nilpotent Collapse]
\lean{ISAR.nilpotent_collapse}
\uses{ISAR.MatrixKernel}
Applying \texttt{konst} to itself maps to the zero matrix, showing the nilpotent collapse of the basis algebra:
$$\text{kernelMatrixView}(\mathbf{K}\ \mathbf{K}) = \mathbf{0}$$
\end{theorem}

\begin{theorem}[Matrix Terminality]
\lean{ISAR.matrix_kernel_terminality}
\uses{ISAR.MatrixKernel, ISAR.ISAR_Kernel_terminal}
Every structure-preserving morphism from \texttt{MatrixKernel} to \texttt{ISAR\_Kernel} is observationally equivalent to the canonical decoding morphism.
\end{theorem}

\begin{theorem}[Unreachability of R and A]
\lean{ISAR.kernel_view_R_unreachable}
\uses{ISAR.MatrixKernel}
The structural homomorphism \texttt{kernelMatrixView} never produces the rotation matrix $R_1$ or adjacency matrix $A_1$ from a pure ISK term.
\end{theorem}

\section{Tensor Denotation Semantics}

We bridge the discrete combinatory logic of the ISAR substrate with continuous analysis and matrix spaces by formalizing the denotational mapping of symbolic ISAR terms into an abstract rank-4 tensor space, where reductions map to extensional equality.

\begin{definition}[Tensor Space Primitives]
\lean{ISAR.ExtTensor}
The abstract tensor space $\mathcal{T}$ is characterized by five primitive operators:
\begin{itemize}
    \item Identity: \texttt{t\_norm}
    \item Constant: \texttt{t\_konst}
    \item Duplicator: \texttt{t\_dup}
    \item Swapper: \texttt{t\_swap}
    \item Composition: \texttt{t\_comp}
\end{itemize}
\end{definition}

\begin{definition}[Derived S in Tensor Space]
\lean{ISAR.t_sₛ}
\uses{ISAR.ExtTensor}
The distributive combinator $\mathbf{S}$ in the tensor space is constructively derived using swapper, composition, and duplicator operators:
\begin{align*}
\text{t\_s}_{\text{s}} &\triangleq \text{t\_app}\Big(\text{t\_app}\big(\text{t\_comp}, \text{t\_app}(\text{t\_comp}, \text{t\_dup})\big), \\
&\quad \text{t\_app}\big(\text{t\_app}(\text{t\_swap}, P), \text{t\_norm}\big)\Big)
\end{align*}
where:
\[ P \triangleq \text{t\_app}\big(\text{t\_app}(\text{t\_comp}, \text{t\_comp}), \text{t\_app}(\text{t\_app}(\text{t\_comp}, \text{t\_comp}), \text{t\_swap})\big) \]
\end{definition}

\begin{definition}[Denotational Mapping]
\lean{ISAR.denot}
\uses{ISAR.ExtTensor, ISAR.t_sₛ}
The structural denotational homomorphism mapping symbolic terms to the tensor space:
$$\text{denot} : \text{ITerm} \to \text{TensorSpace}$$
\end{definition}

\begin{theorem}[Denotational Soundness]
\lean{ISAR.denot_sound}
\uses{ISAR.denot}
One-step reduction in the symbolic calculus maps to extensional equivalence of denotations:
$$\forall t, u, \quad t \to_I u \implies \text{denot } t \approx_{\text{ext}} \text{denot } u$$
Thus, the denotational mapping factors uniquely through the Invariant Layer quotient:
$$\text{InvariantLayer.toExtTensor} : \text{InvariantLayer} \to \text{ExtTensor}$$
\end{theorem}

\subsection{Sparse Matrix Operator Composition}
In the concrete implementation of the axiomatic kernel (\url{scratch/isar_ski_kernel.py}), the emergent operator combinators are represented as products of the basic sparse matrices $I, R, A, S$:
\begin{align*}
\text{NORM} &\triangleq S \cdot I \\
\text{APP} &\triangleq A \cdot R \\
\text{DUP} &\triangleq S \cdot S \\
\text{COMP} &\triangleq R \cdot A \cdot S \\
\text{SWAP} &\triangleq R \cdot S \\
\text{CONST} &\triangleq I \cdot S
\end{align*}

\subsection{Plex-Agent Tensor Primitives and Special Forms}
Alternatively, the runtime environment and compiler stack defined in \url{tensor_primitives.py} specify four basic tensor primitives and four special forms derived from them:
\begin{align*}
\text{ATOM (0)} &\triangleq I \\
\text{PAIR (1)} &\triangleq S \\
\text{JOIN (2)} &\triangleq R \cdot A \\
\text{NORM (3)} &\triangleq I \cdot R \cdot A \cdot S
\end{align*}
The special forms are defined as compositions of these primitives:
\begin{align*}
\text{LAMBDA} &\triangleq \text{JOIN} \cdot \text{PAIR} \cdot \text{NORM} \\
\text{QUOTE} &\triangleq \text{ATOM} \cdot \text{NORM}
\end{align*}
along with conditional branching (\texttt{IF}) and variable bindings (\texttt{DEF}).

\section{Continuous Universal Approximation}

We formalize the continuous parameter space of the ISAR update.

\begin{definition}[Non-Polynomial Activation]
\lean{ISAR.Activation.nonPolynomial}
An activation function $\sigma \in C(\mathbb{R}, \mathbb{R})$ is non-polynomial:
$$\forall P \in \text{Polynomial } \mathbb{R}, \quad \sigma \neq P$$
\end{definition}

\begin{definition}[Raw Address Space]
\lean{ISAR.RawAddress}
The raw parameter configuration space $\text{RawAddress } d\ k$ representing a neural network configuration:
\begin{align*}
\text{RawAddress } d\ k &\triangleq \Sigma (N, T : \mathbb{N}), \quad (\text{Fin } T \to \text{Fin } 4 \to \mathbb{R}) \times C(\mathbb{R}^d, \text{GridState } N) \\
&\quad \times C(\text{GridState } N, \mathbb{R}^k)
\end{align*}
\end{definition}

\begin{definition}[Realization Map]
\lean{ISAR.realizeRaw}
\uses{ISAR.RawAddress}
The continuous function realized by a raw address configuration:
$$\text{realizeRaw } d\ k\ \sigma : \text{RawAddress } d\ k \to C(\mathbb{R}^d, \mathbb{R}^k)$$
\end{definition}

\begin{definition}[Address Equivalence]
\lean{ISAR.AddressEq}
\uses{ISAR.realizeRaw}
Two raw parameter configurations are equivalent if they realize the same continuous function.
\end{definition}

\begin{definition}[Kernel Address Space]
\lean{ISAR.KernelAddress}
\uses{ISAR.AddressEq}
The quotient parameter space of raw addresses modulo observational equivalence:
$$\text{KernelAddress } d\ k\ \sigma := \text{Quotient } (\text{addressSetoid } d\ k\ \sigma)$$
\end{definition}

We state the Universal Approximation Theorem (UAT).

\begin{axiom}[ISAR Universal Approximation]
\lean{ISAR.ISAR_UAT}
\uses{ISAR.Activation.nonPolynomial, ISAR.RawAddress, ISAR.realizeRaw}
Let $K \subset \mathbb{R}^d$ be compact, $f \in C(\mathbb{R}^d, \mathbb{R}^k)$ a target continuous function, $\sigma$ a non-polynomial activation, and $\varepsilon > 0$. Then there exists a raw address $\theta \in \text{RawAddress } d\ k$ such that:
$$\forall x \in K, \quad \| \text{realizeRaw } d\ k\ \sigma\ \theta\ x - f(x) \| < \varepsilon$$
\end{axiom}

\section{Quotient Completion and Unique Representation}

We complete the parameter space and show unique representation.

\begin{axiom}[Quotient Completion Space]
\lean{ISAR.KernelAddressLimit}
The metric completion of \texttt{KernelAddress} is \texttt{KernelAddressLimit}.
\end{axiom}

\begin{axiom}[Continuous Realization Limit]
\lean{ISAR.continuousRealizationLimit}
\uses{ISAR.KernelAddressLimit}
The extended continuous realization map on the completed space.
\end{axiom}

\begin{axiom}[Completion Embedding]
\lean{ISAR.kernelAddressEmbedding}
The canonical inclusion map $i : \text{KernelAddress} \to \text{KernelAddressLimit}$.
\end{axiom}

\begin{theorem}[Embedding Injectivity]
\lean{ISAR.kernelAddressEmbedding_injective}
\uses{ISAR.kernelAddressEmbedding}
The completion embedding is injective:
$$\forall q_1, q_2, \quad i(q_1) = i(q_2) \implies q_1 = q_2$$
\end{theorem}

\begin{theorem}[Embedding Density]
\lean{ISAR.kernelAddressEmbedding_dense}
\uses{ISAR.kernelAddressEmbedding}
The completion embedding has dense range (with respect to uniform convergence on compact domains).
\end{theorem}

\begin{axiom}[Topological Extension Bijection]
\lean{ISAR.topological_extension_bijection}
An injective map with a dense range into a complete metric space extends uniquely to a bijection on the completion of its domain.
\end{axiom}

\begin{theorem}[ISAR Representation Theorem]
\lean{ISAR.ISAR_representation}
\uses{ISAR.topological_extension_bijection, ISAR.kernelAddressEmbedding_injective, ISAR.kernelAddressEmbedding_dense}
Every continuous function $f \in C(\mathbb{R}^d, \mathbb{R}^k)$ has a unique address $\theta_{\text{limit}} \in \text{KernelAddressLimit}$ such that:
$$\text{continuousRealizationLimit } d\ k\ \sigma\ \theta_{\text{limit}} = f$$
\end{theorem}

\begin{theorem}[Physical System Address]
\lean{ISAR.physical_system_has_address}
\uses{ISAR.ISAR_representation}
Every physical system represented as an admissible continuous kernel has a unique address in the continuous completion limit space.
\end{theorem}

\chapter{Discussion, Related Work, and Limitations}

We present a review of related literature, establish the formal boundaries of validity for our theorems, and outline directions for future research.

\section{Related Work}

The ISAR framework intersects several foundational areas of computer science and logic:

\begin{itemize}
    \item \textbf{Combinatory Logic and Lambda Calculus}: The bracket abstraction algorithm compiled to combinator form is a classical approach pioneered by Curry and Feys~\cite{curry1958combinatory} and Barendregt~\cite{barendregt1984lambda}. While traditional implementations suffer from code size explosion (e.g., $O(n^3)$ for standard abstraction), the linear fragment of ISAR maintains bounded sizing, sharing a direct connection to optimal reduction.
    \item \textbf{Optimal Evaluation and Interaction Nets}: Yves Lafont's interaction nets~\cite{lafont1990interaction} provide a graphical formulation of computation where reductions are local and symmetric. Our formalization of the linear fragment (\texttt{LinearIKTerm}) and its bounds mirrors the behavior of optimal virtual machines like the Higher-Order Virtual Machine 2 (HVM2)~\cite{hvm2}, which compiles terms to interaction nets for parallel execution.
    \item \textbf{Categorical Semantics of Computation}: Lawvere's functorial semantics~\cite{lawvere1963functorial} established algebraic theories as categories. Our formulation of semantic kernels ($\mathcal{K}$) and the proof of terminality of the Invariant Layer quotient formalizes view-independence as a universal factorization property.
    \item \textbf{Proof Assistant Kernels}: The verification and automated generation of documentation utilizing Lean 4 quotients and metadata mappings (e.g., doc-gen4~\cite{docgen4}) represent the state-of-the-art in proof design. Our work extends this by proving formal simulation theorems directly against the proof assistant's environment.
\end{itemize}

\section{Limitations and Scope Boundaries}

The theorems presented in this work are accompanied by strict boundary conditions:

\begin{enumerate}
    \item \textbf{Closed Dialects Only}: The terminality and factorization theorems apply exclusively to computational dialects that are \textit{operationally closed}—meaning they possess a total encoding function into the substrate and their evaluation dynamics can be modeled as deterministic, confluent rewrite steps. This excludes open, dynamic, or referentially open systems (such as natural languages, interactive APIs, or non-deterministic transition systems) that require contextual anchors.
    \item \textbf{ZFC HF Set Scope}: The set-theoretic relative interpretation is verified solely for the Hereditarily Finite set fragment. Extension to infinite sets (ZFC with the axiom of infinity) requires different ordinals and is outside the scope of the current Lean 4 proofs.
    \item \textbf{Interpretive Physical Assumptions}: The application of the Universal Approximation Theorem to physical systems rests on the interpretive hypothesis that a physical system's state space can be modeled as an admissible continuous kernel. This hypothesis is a modeling choice and is not mechanically verified by the Lean 4 proof assistant.
\end{enumerate}

\section{Future Work}

This paper forms the core substrate of a multi-part series of monographs:
\begin{itemize}
    \item \textbf{Paper 2 (Continuous Approximation)}: Extension of the discrete kernel to real-valued physical quantities, continuous address spaces, and metric completions, proving the unique continuous representation limit. This builds upon classical universal approximation results on compact domains, particularly those of Cybenko~\cite{cybenko1989approximation} and Leshno et al.~\cite{leshno1993multilayer}.
    \item \textbf{Paper 3 (Hardware Compilation)}: Verification of the compilation path from the algebraic matrix representations to optimal virtual machines (HVM2) and relational database schemas. We also formalize partial evaluation and dialect specialization via the first, second, and third Futamura projections~\cite{futamura1971partial}.
\end{itemize}

\bibliographystyle{plain}
\bibliography{references}



\end{document}
